\section{Problems and Requirements of Typical Usage Scenarios}

Working time in full-body suit with self-contained breathing apparatus is
- usually limited (to two hours at a time OR shorter) and might be reduced even more depending on air and exhaustion levels [\cite{}]
- exhaustion more extreme with more equipment weight + heat [\cite{}]

Additionally, Type 1 suits are not gas-permeable, compromising the body's ability to regulate temperature and leading to even faster exhaustion and heat stress [\cite{holmerProtectiveClothingHot2006}].
Physiological parameters measurable in regular work conditions for catching heat exhaustion early include heart rate (variability), blood pressure [\cite{}].

Other environmental hazards include electricity and possible explosive atmospheres.
- elimination not in task scope, but to be avoided

limited FOV and mobility [\cite{murrayEffectsChemicalProtective2011}]
Work conditions/tasks are characterized by a high mental load [\cite{}].
-both make it hard to keep an eye on all of these things

cognitive load reduced when active?

cognitive offloading

information density?

Task environments can be summarized into two extremes. First, a controlled space with either bright or poor lighting and simple backgrounds. Second, a scenario like an emergency response, in which most variables are unknown. It is likely that backgrounds will be varied and lighting conditions can change.

The types of information needed would differ by task, but life-sustaining paramaters should be front and center in all of them.
